\section{The Ancient Versions and Second Temple Retellings}

Translators who work sentence by sentence through a biblical chapter are its earliest close readers. When a verse carries a chronological snag, ancient translations often smooth it out, insert a clarifying adverb, or sharpen the friction. Laying the Septuagint, the Clementine Vulgate, the Syriac Peshitta, and the Second Temple retellings in \emph{Jubilees} and Josephus beside the Hebrew reveals something unexpected: none of these ancient versions treats verse 5 as a daylight puzzle, yet the Greek Septuagint differs strikingly from the Hebrew and the Vulgate at four verses that bear directly on the chapter's timing and darkness.

\subsection{The Septuagint: Ecstasy, Flame Without Mist, and ``There'' Instead of ``That Day''}

The Greek text of Genesis 15 printed in Henry Barclay Swete's Cambridge edition---where Swete follows Codex Alexandrinus across the opening lacuna of Codex Vaticanus (Genesis 1:1--46:28a)---agrees closely with the Hebrew through the opening six verses.\endnote{Henry Barclay Swete, ed., \emph{The Old Testament in Greek according to the Septuagint}, vol. I, 4th ed. (Cambridge: University Press, 1909), Gen 15:1--21. While Swete's overall base text for Volume I is Codex Vaticanus (B), Vaticanus is missing Genesis 1:1--46:28a in its original fourth-century hand, and Swete supplies those chapters from Codex Alexandrinus (A). For the public-domain English translation of the Sixtine Septuagint text, see Lancelot Charles Lee Brenton, \emph{The Septuagint Version of the Old Testament} (London, 1844).} In verse 1, the word of the Lord comes to Abram \latin{en horamati}, ``in a vision.'' In verse 5, God leads him outside and commands him: \latin{ex\={e}gagen de auton ex\={o} kai eipen pros auton: Anablepson d\={e} eis ton ouranon kai arithm\={e}son tous asteras, ei dyn\={e} exarithm\={e}sai autous}, ``And he brought him out and said to him: Look up now to heaven, and count the stars, if thou shalt be able to number them fully'' (Brenton). The compound verb \latin{anablepson} explicitly adds the upward direction (``look \emph{up}''), and the compound \latin{exarithm\={e}sai} (``count to the end, number fully'') underscores even more clearly than the Hebrew that the obstacle is exhaustive enumeration, not visibility.

From verse 11 onward, however, the Greek diverges from the Masoretic Text at four consecutive joints:

\begin{enumerate}
  \item \textbf{Abram sits beside the carcasses (15:11).} Where the Hebrew vocalizes the verb as \latin{wayyashsh\={e}v} (``and he drove them away''), the Greek translator read the identical consonants as \latin{wayy\={e}shev} (``and he sat'') and rendered: \latin{kai synekathisen autois Abram}, ``and Abram sat down by them.'' That single reading passed through the Old Latin into St. Augustine's \emph{City of God}, where Abram sitting calmly beside the divided birds becomes an image of the faithful persevering amid carnal divisions.
  \item \textbf{Deep sleep becomes ecstasy (15:12).} For Hebrew \latin{wayeh\={\i} ha-shemesh l\={a}v\={o}' we-tard\={e}m\={a}h n\={a}fel\={a}h}, the Greek translates \latin{peri de h\={e}liou dysmas ekstasis epepesen t\={o} Abram, kai idou phobos skotinos megas epipiptei aut\={o}}, ``And about sunset a trance fell upon Abram, and lo! a great gloomy terror falls upon him.'' By turning \latin{tard\={e}m\={a}h} (deep sleep) into \latin{ekstasis} (trance or rapture out of one's senses), the Septuagint pulls verse 12 toward visionary transport rather than ordinary bodily fatigue.
  \item \textbf{The sun still setting, and flame in place of dark mist (15:17).} Swete's Greek reads \latin{epei de egineto ho h\={e}lios pros dysmais, phlox egeneto; kai idou klibanos kapnizomenos kai lampades pyros}, ``And when the sun was about to set, there was a flame, and behold a smoking furnace and lamps of fire.'' Two features stand out immediately. First, the imperfect \latin{egineto ... pros dysmais} (``was drawing toward setting'') leaves the sun still going down rather than already set, matching a participial reading of Hebrew \latin{b\={a}'\={a}h}. Second, the Greek carries \emph{no equivalent whatsoever} for Hebrew \latin{'al\={a}\d{t}\={a}h} (``thick darkness'' or ``dark mist''): in its place stands \latin{phlox egeneto}, ``there was a flame,'' followed by plural \latin{lampades pyros}, ``lamps of fire.''
  \item \textbf{``There'' instead of ``on that day'' (15:18).} At the opening of verse 18, Swete's Alexandrinus text reads \latin{ekei dietheto ho Theos t\={o} Abram diath\={e}k\={e}n}, ``\emph{There} God made a covenant with Abram,'' replacing the temporal phrase ``on that day'' (\latin{ba-yy\={o}m ha-h\={u}'}) with a spatial adverb! By contrast, the Sixtine Septuagint text translated by Brenton---following the fifteenth-century minuscule supplement supplied to Codex Vaticanus---and the Old Latin cited by St. Augustine preserve \latin{en t\={e} h\={e}mera ekein\={e}}, ``In that day.''
\end{enumerate}

The net effect of the Septuagint is remarkable. While the Greek keeps sunset after the star-promise, it weakens the daytime-to-black-night contrast at verses 17 and 18 (replacing thick darkness with flame and, in Codex Alexandrinus, replacing ``on that day'' with ``there''), and by translating \latin{tard\={e}m\={a}h} as \latin{ekstasis} it aligns the whole scene with prophetic trance.

\subsection{The Vulgate and the Syriac Peshitta}

St. Jerome's Latin translation from the Hebrew, preserved in the Clementine Vulgate and mirrored in the Douay-Rheims Bible, restores the Masoretic profile with precision. In verse 1, the word comes \latin{per visionem} (``by a vision''); in verse 5, God leads Abram out (\latin{Eduxitque eum foras}) and commands \latin{Suspice c\ae lum, et numera stellas, si potes} (``Look up to heaven and number the stars if thou canst''); in verse 11, Jerome corrects the Septuagint's sitting into active guarding (\latin{abigebat eas Abram}, ``Abram drove them away''); and in verse 12, he replaces \latin{ekstasis} with \latin{sopor} (``a deep sleep'').

Crucially, Jerome heightens the solar progression between verses 12 and 17 through his choice of Latin tenses. Verse 12 uses the imperfect subjunctive, \latin{Cumque sol occumberet} (``And when the sun was setting''), whereas verse 17 switches to the pluperfect subjunctive, \latin{Cum ergo occubuisset sol, facta est caligo tenebrosa} (``And when the sun was set, there arose a dark mist''), followed by \latin{In illo die} (``That day'') in verse 18. Ironically, therefore, the Latin Vulgate and its English heir, the Douay-Rheims, make the march from afternoon toward sunset (v. 12) and full nightfall (v. 17) even crisper to a Western reader than the Septuagint did---yet Jerome himself, in his \emph{Hebrew Questions on Genesis}, saw no clash at all with verse 5.

The Syriac Peshitta stands squarely beside the Hebrew across all four temporal joints.\endnote{Syriac Peshitta on Genesis 15:5, 12, and 17--18.} It preserves God bringing Abram outside to look toward heaven and number the stars at 15:5, the sun approaching setting at 15:12 and again at 15:17, and ``on that day'' at 15:18, inserting no explicit elapsed interval between verses 5 and 12.

\subsection{Second Temple Retellings: Jubilees and Josephus}

Two Jewish writings from the Second Temple period directly inspected here retell Genesis 15 in strikingly different ways.

The \emph{Book of Jubilees} (second century BC) rewrites Genesis 15 in chapter 14:1--20 with meticulous attention to its own liturgical calendar.\endnote{R. H. Charles, trans., \emph{The Book of Jubilees}, in \emph{The Apocrypha and Pseudepigrapha of the Old Testament}, vol.~II (Oxford: Clarendon Press, 1913), ch.~14:1--20.} At 14:5, where Genesis says only that God commanded Abram to look and count, Jubilees explicitly adds that Abram looked up and \emph{beheld the stars}. Five verses later (14:10), Jubilees dates the covenant sacrifice to the middle of the third month---the festival of Weeks (Shavuot) in Jubilees' solar calendar---then reports sunset completed at 14:13 and 14:17, and retains ``on that day'' at 14:18. For a work obsessed with exact dates and calendrical consistency, Jubilees assumes without a flicker of anxiety that Abram actually saw real stars in verse 5, while inserting no explicit day-break before verse 13.

Flavius Josephus (AD~37--c.~100) takes the opposite path in \emph{Jewish Antiquities} I.10.3.\endnote{Flavius Josephus, \emph{Antiquitates Judaicae} I.10.3 (trans. William Whiston, London, 1737).} Compressing Genesis 15 into a brisk historical summary for Greco-Roman readers, Josephus keeps God's promise that Abram's posterity would rival the multitude of the stars and passes immediately to the commanded sacrifice, omitting altogether the stepping outside, the command to look up, the setting of the sun, the darkness, and the concluding formula ``on that day.''

Across every inspected ancient biblical version and Second Temple retelling, the result is unanimous. Translators and retellers noticed the birds of verse 11, the trance of verse 12, the darkness or flame of verse 17, and the covenant date of Jubilees 14:10. Not one of them thought Genesis 15 depicted Abram counting invisible stars in daylight.
